\documentclass[../preamble]{subfiles}
\begin{document}

\section{Product Category}
\label{sec:product-category}
\concept{product-category}{Product Category}
\depends{category}
\depends{functor}
\implements{haskell}{Cat.Product}
\implements{agda}{Cat.Product}
\implements{lisp}{product-category}

Given categories $\Category{C}$ and $\Category{D}$,
their \newterm{product category}
$\newmath{\Category{C} \times \Category{D}}$
is defined as follows
(\nlabref{product+category}):

\begin{definition}[Product category]
  \label{def:product-category}
  The \newterm{product category} $\Category{C} \times \Category{D}$
  consists of:
  \begin{enumerate}
    \item \textbf{Objects.}
      $\Ob(\Category{C} \times \Category{D})
        = \Ob(\Category{C}) \times \Ob(\Category{D})$.
      An object is a pair $\newmath{(X, Y)}$
      with $X \in \Ob(\Category{C})$ and $Y \in \Ob(\Category{D})$.

    \item \textbf{Morphisms.}
      A morphism $(X_1, Y_1) \lto (X_2, Y_2)$
      in $\Category{C} \times \Category{D}$
      is a pair $\newmath{(f, g)}$
      with $f \colon X_1 \lto X_2$ in $\Category{C}$
      and $g \colon Y_1 \lto Y_2$ in $\Category{D}$.

    \item \textbf{Identity.}
      $\newmath{\id_{(X,Y)} = (\id_X, \id_Y)}$.

    \item \textbf{Composition.}
      $\newmath{(f_2, g_2) \circ (f_1, g_1) = (f_2 \circ f_1, g_2 \circ g_1)}$.
  \end{enumerate}
\end{definition}

\begin{proposition}
  \label{prop:product-cat-laws}
  The product category $\Category{C} \times \Category{D}$
  satisfies the category axioms.
\end{proposition}

\begin{proof}
  Associativity and identity laws hold componentwise.
  For associativity:
  \begin{align*}
    (h_1, h_2) \circ ((g_1, g_2) \circ (f_1, f_2))
    &= (h_1, h_2) \circ (g_1 \circ f_1, g_2 \circ f_2) \\
    &= (h_1 \circ (g_1 \circ f_1), h_2 \circ (g_2 \circ f_2)) \\
    &= ((h_1 \circ g_1) \circ f_1, (h_2 \circ g_2) \circ f_2) \\
    &= (h_1 \circ g_1, h_2 \circ g_2) \circ (f_1, f_2) \\
    &= ((h_1, h_2) \circ (g_1, g_2)) \circ (f_1, f_2),
  \end{align*}
  using associativity in $\Category{C}$ and $\Category{D}$ respectively.
  Left and right identity follow similarly from the component identities.
\end{proof}

\subsection{Projection Functors}

\begin{definition}[Projection functors]
  \label{def:projection-functors}
  The \newterm{projection functors} are:
  \begin{align}
    \newmath{\Functor{\pi}_1} &\colon \Category{C} \times \Category{D} \lto \Category{C},
      &(X, Y) &\lmapsto X,
      &(f, g) &\lmapsto f, \label{eq:pi1} \\
    \newmath{\Functor{\pi}_2} &\colon \Category{C} \times \Category{D} \lto \Category{D},
      &(X, Y) &\lmapsto Y,
      &(f, g) &\lmapsto g. \label{eq:pi2}
  \end{align}
  These are strict functors:
  they preserve identities and composition by definition.
\end{definition}

\subsection{Universal Property}

\begin{proposition}[Universal property of the product]
  \label{prop:product-universal}
  For any category $\Category{E}$
  and functors $\Functor{F} \colon \Category{E} \lto \Category{C}$
  and $\Functor{G} \colon \Category{E} \lto \Category{D}$,
  there exists a unique functor
  $\newmath{\langle \Functor{F}, \Functor{G} \rangle}
    \colon \Category{E} \lto \Category{C} \times \Category{D}$
  such that
  $\Functor{\pi}_1 \circ \langle \Functor{F}, \Functor{G} \rangle = \Functor{F}$
  and
  $\Functor{\pi}_2 \circ \langle \Functor{F}, \Functor{G} \rangle = \Functor{G}$.
\end{proposition}

This universal property is expressed by the following diagram:
\begin{equation}
  \label{eq:product-universal-diagram}
  \begin{tikzcd}
    & \Category{E}
      \arrow[dl, "\Functor{F}"']
      \arrow[dr, "\Functor{G}"]
      \arrow[d, dashed, "{\langle \Functor{F}{,} \Functor{G} \rangle}"] \\
    \Category{C}
    & \Category{C} \times \Category{D}
      \arrow[l, "\Functor{\pi}_1"]
      \arrow[r, "\Functor{\pi}_2"']
    & \Category{D}
  \end{tikzcd}
\end{equation}

\begin{proof}
  Define $\langle \Functor{F}, \Functor{G} \rangle$ on objects by
  $E \lmapsto (\Functor{F}(E), \Functor{G}(E))$
  and on morphisms by
  $h \lmapsto (\Functor{F}(h), \Functor{G}(h))$.
  Functoriality follows from that of $\Functor{F}$ and $\Functor{G}$.
  The projection equations hold by definition.
  Uniqueness: if $\Functor{H}$ satisfies both equations,
  then on objects $\Functor{H}(E)
    = (\Functor{\pi}_1(\Functor{H}(E)), \Functor{\pi}_2(\Functor{H}(E)))
    = (\Functor{F}(E), \Functor{G}(E))$,
  and similarly on morphisms.
\end{proof}

\end{document}
